\chapter{Gastroscopy}

\section*{Definition and Overview}
A gastroscopy, also called an upper endoscopy or oesophago-gastro-duodenoscopy (OGD), is a procedure in which a thin, flexible tube with a camera is passed through the mouth into the oesophagus, stomach, and the first part of the small intestine. It allows the doctor to see the lining of these organs directly, take biopsies, and treat some conditions during the same procedure. Gastroscopy is used to investigate symptoms such as heartburn, difficulty swallowing, abdominal pain, vomiting, and bleeding, and to diagnose ulcers, inflammation, and cancers. It is performed under sedation, takes about 10--15 minutes, and most people go home the same day. This chapter explains why gastroscopy is done, what it involves, and what to expect.

\section*{Epidemiology}
Gastroscopy is one of the most commonly performed endoscopic procedures in Australia, with hundreds of thousands performed each year. It is the standard investigation for upper gastrointestinal symptoms and is essential in the diagnosis of ulcers, gastritis, coeliac disease, and oesophageal and stomach cancer. Around 1 in 5 Australians has reflux, and gastroscopy is used selectively for warning symptoms, persistent symptoms, or screening for Barrett's oesophagus. The procedure is safe, with serious complications occurring in well under 1% of cases, and it is performed as a day procedure.

\section*{Aetiology and Pathophysiology}
The procedure diagnoses and treats conditions of the upper gut.
\begin{itemize}
    \item \textbf{Visualisation:} the endoscope with a camera and light allows direct inspection of the oesophagus, stomach, and duodenum
    \item \textbf{Diagnosis:} oesophagitis, gastritis, ulcers, coeliac disease, Barrett's oesophagus, and cancer are seen directly
    \item \textbf{Biopsy:} small tissue samples are taken painlessly through the scope for pathology
    \item \textbf{H.\ pylori testing:} biopsy samples detect the bacterium that causes ulcers and gastritis
    \item \textbf{Treatment:} bleeding ulcers can be injected or clipped, polyps removed, and strictures dilated during the procedure
    \item \textbf{Sedation:} a light sedative and throat spray make the procedure comfortable
    \item \textbf{The scope:} about the width of a finger, passed carefully through the mouth
    \item \textbf{Why it matters:} direct visualisation and biopsy are far more accurate than imaging alone for upper gut disease
\end{itemize}

\section*{Clinical Features}
\begin{itemize}
    \item The procedure takes 10--15 minutes under light sedation
    \item A local anaesthetic throat spray numbs the throat
    \item You lie on your left side, and the scope is passed gently into the oesophagus
    \item Air is used to open the stomach for a clear view, causing some fullness
    \item Biopsies and treatments do not cause pain
    \item Afterwards: drowsiness from sedation, and a mildly sore throat for a day
    \item You cannot drive for 24 hours after sedation
    \item Results are usually given immediately, with biopsy results days later
\end{itemize}

\section*{Differential Diagnosis}
Gastroscopy is performed to distinguish the causes of upper gut symptoms.
\begin{itemize}
    \item Gastro-oesophageal reflux disease and oesophagitis
    \item Gastritis and peptic ulcers, with H.\ pylori
    \item Coeliac disease, confirmed on biopsy
    \item Barrett's oesophagus and oesophageal or stomach cancer
    \item Strictures causing difficulty swallowing
    \item Functional dyspepsia, where the endoscopy is normal
    \item The findings guide treatment
\end{itemize}

\section*{When to Seek Medical Care}
Gastroscopy is arranged by a doctor for specific indications, including warning symptoms such as difficulty swallowing, weight loss, vomiting blood, or black stools, and for persistent symptoms not responding to treatment.

\textbf{Call 000 (or local emergency services) for:}
\begin{itemize}
    \item Severe abdominal pain, fever, or vomiting blood after the procedure
    \item Difficulty breathing or chest pain after sedation
    \item Heavy bleeding from the back passage
    \item Signs of a perforation: severe pain, fever, and a rigid abdomen
    \item Any emergency symptoms
\end{itemize}

\section*{Diagnosis and Investigations}
\begin{itemize}
    \item \textbf{Pre-procedure preparation:} fasting for 6 hours before the procedure, and advice about medicines, especially blood thinners
    \item \textbf{The endoscopy:} the diagnostic procedure, with biopsy and treatment as needed
    \item \textbf{Histology:} pathology of biopsies, which is essential for coeliac disease, H.\ pylori, and cancer
    \item \textbf{Blood tests:} for anaemia and other findings as indicated
    \item \textbf{Imaging:} occasionally CT for staging or further assessment
\end{itemize}

\section*{Management}
\begin{itemize}
    \item \textbf{Before:} fasting, and stopping some medicines on medical advice
    \item \textbf{During:} sedation and throat spray, with monitoring by the team
    \item \textbf{After:} rest until the sedation wears off, then light food and drink
    \item \textbf{Treatment of findings:} acid suppression for reflux and oesophagitis, antibiotics for H.\ pylori, gluten-free diet for coeliac disease, and specialist treatment for cancer and Barrett's oesophagus
    \item \textbf{Strictures:} dilatation during the procedure
    \item \textbf{Bleeding:} endoscopic treatment at the time of diagnosis
    \item \textbf{Follow-up:} repeat endoscopy for surveillance of Barrett's oesophagus, ulcers, and after treatment
\end{itemize}

\section*{Complications}
\begin{itemize}
    \item Sore throat, which is common and minor
    \item Reactions to sedation, including breathing and heart effects, which are rare
    \item Bleeding after biopsy or polyp removal
    \item Perforation of the oesophagus or stomach, a rare but serious complication
    \item Aspiration of stomach contents during the procedure, rare with fasting
    \item The risks of not investigating: missed diagnoses of serious conditions
\end{itemize}

\section*{Prognosis and Outcomes}
Gastroscopy is safe, quick, and highly accurate, and it is the definitive investigation for upper gut disease. Most people have the procedure and return to normal life the same day, and the findings guide effective treatment: ulcers heal, H.\ pylori is eradicated, coeliac disease is managed with diet, and cancers are detected early, when they are most treatable. The procedure's value in diagnosis and its ability to treat during the same visit make it one of the most useful tools in digestive medicine.

\section*{Prevention and Self-Care}
\begin{itemize}
    \item Discuss the reason for the procedure and the risks with the doctor before consenting
    \item Follow the fasting instructions exactly
    \item Tell the doctor about all medicines, including blood thinners, and allergies
    \item Arrange for someone to drive you home after sedation
    \item Do not drive, operate machinery, or make important decisions for 24 hours
    \item Rest and eat lightly after the procedure
    \item Report severe pain, fever, or bleeding promptly
    \item Attend follow-up and surveillance procedures as recommended
\end{itemize}

\section*{Sources and Further Reading}
The following reputable sources provide further information for healthcare students and patients seeking reliable, current advice about gastroscopy.
\begin{itemize}
    \item Healthdirect Australia. \textit{Gastroscopy.} https://www.healthdirect.gov.au/gastroscopy
    \item Gastroenterology Society of Australia (GESA). \textit{Gastroscopy patient information.} https://www.gesa.org.au/
    \item Better Health Channel. \textit{Gastroscopy.} https://www.betterhealth.vic.gov.au/
    \item NHS. \textit{Gastroscopy.} https://www.nhs.uk/conditions/gastroscopy/
    \item Mayo Clinic. \textit{Upper endoscopy.} https://www.mayoclinic.org/
    \item MedlinePlus. \textit{Upper GI endoscopy.} https://medlineplus.gov/
\end{itemize}
